\documentclass[12pt,a4paper]{article}

% ----------------------------------------------------------- Pakete
\usepackage[ngerman]{babel}
\usepackage[T1]{fontenc}
\usepackage[utf8]{inputenc}
\usepackage{lmodern}
\usepackage{microtype}
\usepackage{xcolor}
\usepackage{geometry}
\geometry{margin=2.5cm}
\usepackage{amsmath,amssymb,mathtools}
\usepackage{graphicx}
\usepackage{booktabs}
\usepackage{siunitx}
\usepackage{hyperref}
\hypersetup{colorlinks=true,linkcolor=blue!50!black,citecolor=green!40!black,urlcolor=blue!60!black}
\usepackage{cleveref}

\title{Ein modernes \LaTeX-Dokument}
\author{Max Mustermann \\ \texttt{max@example.org}}
\date{\today}

\begin{document}
\maketitle

\begin{abstract}
  Dieses Dokument demonstriert die wichtigsten Funktionen eines vollständigen
  \LaTeX-Editors: Syntaxhervorhebung, Autocomplete, Live-Vorschau, SyncTeX,
  Gliederungs-Outline und Fehleranalyse.
\end{abstract}

\tableofcontents
\newpage

\section{Einleitung}
\label{sec:intro}
\LaTeX\ ist das Standard-System für wissenschaftliche Publikationen. Es trennt
Inhalt und Formatierung zuverlässig voneinander~\cite{lamport1994}.

Ein Referenz auf \cref{sec:math} zeigt, wie \texttt{cleveref} automatisch
nummerierte Verweise erzeugt.

\section{Mathematik}
\label{sec:math}

Die Euler-Identität ist eine der schönsten Gleichungen der Mathematik:
\begin{equation}
  \label{eq:euler}
  e^{i\pi} + 1 = 0
\end{equation}

Vektoren und Matrizen lassen sich bequem setzen:
\begin{align}
  \vec{a} \times \vec{b} &= \begin{pmatrix} a_2b_3 - a_3b_2 \\ a_3b_1 - a_1b_3 \\ a_1b_2 - a_2b_1 \end{pmatrix}
  \label{eq:cross}\\
  A^{-1} &= \frac{1}{\det(A)}\,\operatorname{adj}(A)
  \label{eq:inv}
\end{align}

Nach \cref{eq:euler} und \cref{eq:cross} folgt ein weiteres Beispiel mit
SI-Einheiten: \SI{9{,}81}{\meter\per\square\second}.

\section{Listen und Tabellen}

\subsection{Aufzählung}
\begin{itemize}
  \item Vorteil: hohe Typografiequalität
  \item Vorteil: Trennung von Inhalt und Gestaltung
  \begin{enumerate}
    \item verschachtelte Nummerierung
    \item weitere Ebene
  \end{enumerate}
\end{itemize}

\subsection{Tabelle}
\label{tab:vergleich}
\begin{table}[htbp]
  \centering
  \caption{Vergleich von Satzsystemen}
  \label{tab:satzsysteme}
  \begin{tabular}{lrr}
    \toprule
    System        & {Jahrgang} & {Marktanteil (\%)} \\
    \midrule
    \LaTeX        & 1984 & 31,2 \\
    Word          & 1983 & 55,7 \\
    Troff/nroff   & 1973 &  2,1 \\
    \bottomrule
  \end{tabular}
\end{table}

Siehe \cref{tab:vergleich} und \cref{tab:satzsysteme}.

\section{Abbildungen}
\label{sec:figures}

\begin{figure}[htbp]
  \centering
  \fbox{\parbox{0.6\linewidth}{\centering\vspace{1.5cm}
    Platzhalter für eine Grafik\\
    (\texttt{includegraphics} einsetzen)\vspace{1.5cm}}}
  \caption{So sieht eine Abbildung aus.}
  \label{fig:placeholder}
\end{figure}

\section{Literatur}
\label{sec:literatur}

Zitiert wird mit \texttt{\textbackslash cite}\texttt{\{...\}}, siehe
\cite{lamport1994,kopka2003,mittelbach2004}. Das Literaturverzeichnis wird
automatisch aus der Datei \texttt{references.bib} erzeugt.

\section{Fazit}
\label{sec:fazit}
Alles funktioniert – von \cref{sec:intro} bis hierhin.

% TODO: Zusammenfassung noch ergänzen

\newpage
\bibliographystyle{plain}
\bibliography{references}

\end{document}
